\documentclass{umk-matematika}

\begin{document}

\maketitlepage
    {Практическая работа №20}
    {Параллелепипед. Объём}
    {}

\section{Фундамент (1--4 балла)}

\textbf{Задача 1 (1 балл).} Измерения прямоугольного параллелепипеда: $3$ см, $4$ см, $5$ см. Найдите его объём.

\textbf{Задача 2 (2 балла).} Объём прямоугольного параллелепипеда равен $120$ см³. Два измерения равны $4$ см и $5$ см. Найдите третье измерение.

\textbf{Задача 3 (3 балла).} Куб имеет ребро $4$ см. Найдите его объём.

\textbf{Задача 4 (4 балла).} Площадь основания прямоугольного параллелепипеда $15$ см², высота $7$ см. Найдите объём.

\section{Стены (5--7 баллов)}

\textbf{Задача 5 (5 баллов).} Диагональ прямоугольного параллелепипеда равна $13$ см. Измерения основания $3$ см и $4$ см. Найдите объём.

\textbf{Задача 6 (6 баллов).} Объём куба равен $125$ см³. Найдите площадь его полной поверхности.

\textbf{Задача 7 (7 баллов).} В прямоугольном параллелепипеде стороны основания $6$ см и $8$ см. Диагональ параллелепипеда наклонена к плоскости основания под углом $45^\circ$. Найдите объём.

\section{Кровля (8--10 баллов)}

\textbf{Задача 8 (8 баллов).} Фундамент здания — прямоугольный параллелепипед с размерами $12$ м × $8$ м × $1{,}5$ м. Внутри вырезан прямоугольный проём для подвала размером $6$ м × $4$ м на всю высоту фундамента. Найдите объём бетона.

\textbf{Задача 9 (9 баллов).} Бак для воды имеет форму прямоугольного параллелепипеда с квадратным основанием. Объём бака $1000$ литров, высота $1$ м. Найдите сторону основания в метрах. (1 м³ = 1000 л)

\textbf{Задача 10 (10 баллов).} Комната имеет форму прямоугольного параллелепипеда. Площадь пола $20$ м², площадь одной стены $15$ м², площадь смежной стены $12$ м². Найдите объём комнаты.

\newpage
\section*{Ответы}

\begin{enumerate}
  \item $60$ см³
  \item $6$ см
  \item $64$ см³
  \item $105$ см³
  \item $144$ см³
  \item $150$ см²
  \item $480$ см³
  \item $108$ м³
  \item $1$ м
  \item $60$ м³
\end{enumerate}

\end{document}